\subsection{Abstand eines Punktes von einer Geraden}
\label{subsec:analytic_geometry_point_to_line_distance}

Die kürzeste Verbindung eines Punktes mit einer Geraden steht senkrecht auf
der Geraden. Sie entsteht, indem der Anteil in Geradenrichtung vom
Verbindungsvektor abgezogen wird.

\begin{center}
  \CachedTikzFigure{content/resources/chapters/analytic_geometry/tikz/point_line_distance.tex}
\end{center}

\begin{TheoremBox}{Punkt--Geraden-Abstand}{analytic_geometry_point_line_distance}
  Gegeben seien
  \[
    g:\quad\vec x=\vec p+\lambda\vec a,
    \qquad Q\text{ mit Ortsvektor }\vec q.
  \]
  Mit $\vec r=\vec q-\vec p$ und
  $\vec e_a=\vec a/\lvert\vec a\rvert$ gilt
  \[
    \vec r_{\parallel}=(\vec r\cdot\vec e_a)\vec e_a,
    \qquad
    d(Q,g)=\left\lvert\vec r-\vec r_{\parallel}\right\rvert.
  \]
\end{TheoremBox}

\begin{ExampleBox}{Abstand eines Punktes von einer Geraden}{analytic_geometry_point_line_distance_example}
  Gegeben sind
  \[
    Q(0,2),
    \qquad
    g:\quad\vec x=\begin{pmatrix}2\\2\end{pmatrix}
      +\lambda\begin{pmatrix}1\\-1\end{pmatrix}.
  \]
  Wir zerlegen den Verbindungsvektor in einen parallelen und einen
  senkrechten Anteil:
  \begin{ClozeArea}[34]{16}
    \[
      \begin{aligned}
        \vec r
          &=\begin{pmatrix}0\\2\end{pmatrix}
            -\begin{pmatrix}2\\2\end{pmatrix}
            =\begin{pmatrix}-2\\0\end{pmatrix},
        &
        \vec e_a
          &=\frac1{\sqrt2}\begin{pmatrix}1\\-1\end{pmatrix},\\[1em]
        \vec r_{\parallel}
          &=(\vec r\cdot\vec e_a)\vec e_a
            =-\sqrt2\,\vec e_a
            =\begin{pmatrix}-1\\1\end{pmatrix},
        &
        \vec r_{\perp}
          &=\vec r-\vec r_{\parallel}
            =\begin{pmatrix}-1\\-1\end{pmatrix}.
      \end{aligned}
    \]
  \end{ClozeArea}
  Der gesuchte Abstand ist daher
  \[
    d(Q,g)=\Cloze{7cm}{\lvert\vec r_{\perp}\rvert=\sqrt2}.
  \]
\end{ExampleBox}
